% Provenance (never printed): see sources/outline_sources.md and sources/book1/BRIEFS.md (g6/04).
\chapter{Decisions: \texttt{if} and \texttt{else}}\label{ch:g6:if-else}

At the cinema, a ticket costs 5 euros for a child under 12, 6 euros from the
age of 65 and 9 euros for everyone else. The ticket machine asks only one
question, ``How old are you?'', and then decides which price to print. It does
not run the same \omterm{def:g1:instructions-in-order:instruction}{instructions} for everybody: it chooses. In this chapter,
our \omterm{def:g4:programs-and-operating-system:program}{programs} learn to choose.

\begin{recall}
A \omterm{def:g4:algorithms-repeat-and-decide:condition}{condition} is a question whose answer is yes or no, and an \omterm{def:g4:algorithms-repeat-and-decide:algorithm}{algorithm} can do
one thing or another depending on it (\cref{ch:g4:algorithms-repeat-and-decide}).
\code{input} reads a typed answer as text, and \code{int} converts it into a
whole number (\cref{ch:g6:variables-and-input}). A \code{for} \omterm{def:g4:algorithms-repeat-and-decide:loop}{loop} repeats an
indented block (\cref{ch:g6:for-loops-turtle}).
\end{recall}

\section{Comparisons give true or false}

\begin{definition}[Comparison and Boolean value]\label{def:g6:if-else:comparison}
A \emph{comparison}\index{comparison} asks how two \omterm{def:g6:variables-and-input:value}{values} compare, with one
of the signs \code{==} (equal), \code{!=} (not equal), \code{<}, \code{<=},
\code{>} and \code{>=}. Its result is a \emph{Boolean value}\index{Boolean
value}: \code{True} or \code{False}, the only two \omterm{def:g6:variables-and-input:value}{values} of the type
\code{bool}.
\end{definition}

\Needspace*{9\baselineskip}
\omcode{code/grade-6/04-if-else/listings/compare.py}{1}{6}{Six
\omterm{def:g6:if-else:comparison}{comparisons}.}\label{lst:g6:if-else:compare}

\Needspace*{10\baselineskip}
\omoutput{code/grade-6/04-if-else/outputs/compare.txt}{Capital and small letters
are different, so \code{"cat"} is not equal to
\code{"Cat"}.}\label{out:g6:if-else:compare}

Be careful with the two signs: one \code{=} puts a \omterm{def:g6:variables-and-input:value}{value} in a \omterm{def:g6:variables-and-input:value}{variable}, two
\code{==} ask whether two \omterm{def:g6:variables-and-input:value}{values} are equal. Writing one instead of two in a
question is a common mistake, and Python refuses it:

\Needspace*{9\baselineskip}
\omoutput{code/grade-6/04-if-else/outputs/assign_run.txt}{\code{if x = 3:}
is not a question. Python even suggests the
fix.}\label{out:g6:if-else:assign}

\section{\texorpdfstring{\texttt{if}}{if} and \texorpdfstring{\texttt{else}}{else}}

\begin{definition}[Conditional statement and branch]\label{def:g6:if-else:if}
A \emph{conditional statement}\index{conditional statement}, or \emph{if
statement}\index{if statement}, runs a block only when a \omterm{def:g4:algorithms-repeat-and-decide:condition}{condition} is true:
\code{if condition:} followed by an indented block. It may have an
\code{else:} followed by a second block, run when the \omterm{def:g4:algorithms-repeat-and-decide:condition}{condition} is false. Each
block is a \emph{branch}\index{branch}: on each run, exactly one branch is
taken.
\end{definition}

\Needspace*{9\baselineskip}
\omcode{code/grade-6/04-if-else/listings/even_odd.py}{1}{6}{Even or
odd?}\label{lst:g6:if-else:even-odd}

\Needspace*{7\baselineskip}
\omoutput{code/grade-6/04-if-else/outputs/even_odd.txt}{17 divided by 2 leaves
1, so the \code{else} \omterm{def:g6:if-else:if}{branch} ran; line 6 runs in every
case.}\label{out:g6:if-else:even-odd}

As for the \code{for} \omterm{def:g4:algorithms-repeat-and-decide:loop}{loop}, the line before a block ends with a colon, and the
block is indented. Line 6 is not indented: it is after the \code{if} statement
and runs whichever \omterm{def:g6:if-else:if}{branch} was taken.

\begin{omfigure}
\begin{tikzpicture}[x=1cm, y=1cm, font=\footnotesize,
    q/.style={draw=black!70, diamond, aspect=2.4, fill=ocYellow!35, inner sep=1pt, align=center},
    blk/.style={draw=black!65, rounded corners=3pt, fill=#1, minimum width=2.9cm, minimum height=0.8cm, align=center}]
  \node[blk=black!6] (in) at (0,2.2) {\code{n = int(input(...))}};
  \node[q] (c) at (0,0.8) {\texttt{n \% 2 == 0}?};
  \node[blk=ocMint!20] (yes) at (-2.6,-0.7) {\code{print(n, "is even")}};
  \node[blk=ocRed!12] (no) at (2.6,-0.7) {\code{print(n, "is odd")}};
  \node[blk=black!6] (end) at (0,-2.2) {\code{print("Thank you!")}};
  \draw[-{Stealth[length=2mm]}, thick] (in) -- (c);
  \draw[-{Stealth[length=2mm]}, thick] (c.west) -| node[pos=0.25, above] {\code{True}} (yes.north);
  \draw[-{Stealth[length=2mm]}, thick] (c.east) -| node[pos=0.25, above] {\code{False}} (no.north);
  \draw[-{Stealth[length=2mm]}, thick] (yes.south) |- (end.west);
  \draw[-{Stealth[length=2mm]}, thick] (no.south) |- (end.east);
\end{tikzpicture}

{\small An \code{if} statement is a fork: the \omterm{def:g4:algorithms-repeat-and-decide:condition}{condition} sends the \omterm{def:g4:programs-and-operating-system:program}{program} into
one of the two blocks, and both roads meet again after them.}
\end{omfigure}

\section{\texorpdfstring{\texttt{elif}}{elif}: more than two cases}

The cinema has three prices. Between \code{if} and \code{else}, any number of
\code{elif} parts (short for ``else if'') can be added, each with its own
\omterm{def:g4:algorithms-repeat-and-decide:condition}{condition}. Python checks the \omterm{def:g4:algorithms-repeat-and-decide:condition}{conditions} in order, from top to bottom, and runs
the block of the first one that is true; it skips all the others. The
\code{else} block runs only if no \omterm{def:g4:algorithms-repeat-and-decide:condition}{condition} was true.

\Needspace*{11\baselineskip}
\omcode{code/grade-6/04-if-else/listings/ticket.py}{1}{8}{The ticket
machine.}\label{lst:g6:if-else:ticket}

\Needspace*{6\baselineskip}
\omoutput{code/grade-6/04-if-else/outputs/ticket.txt}{70 is not less than 12,
nor less than 65: the \code{else} \omterm{def:g6:if-else:if}{branch}.}\label{out:g6:if-else:ticket}

\begin{omfigure}
\begin{tikzpicture}[x=1cm, y=1cm, font=\footnotesize,
    q/.style={draw=black!70, rounded corners=6pt, fill=ocYellow!35, minimum width=2.6cm, minimum height=0.75cm},
    ans/.style={draw=black!65, rounded corners=2pt, fill=ocBlue!10, minimum width=2.1cm, minimum height=0.7cm}]
  \node[q] (q1) at (0,0) {\code{age < 12}?};
  \node[q] (q2) at (0,-1.4) {\code{age < 65}?};
  \node[ans] (a1) at (4,0) {price 5};
  \node[ans] (a2) at (4,-1.4) {price 9};
  \node[ans] (a3) at (0,-2.8) {price 6};
  \draw[-{Stealth[length=2mm]}, thick] (q1) -- node[above] {yes} (a1);
  \draw[-{Stealth[length=2mm]}, thick] (q2) -- node[above] {yes} (a2);
  \draw[-{Stealth[length=2mm]}, thick] (q1) -- node[right] {no} (q2);
  \draw[-{Stealth[length=2mm]}, thick] (q2) -- node[right] {no, \code{else}} (a3);
  \draw[-{Stealth[length=2mm]}, thick] (0,0.9) -- (q1);
\end{tikzpicture}

{\small The \code{elif} chain is a ladder of questions, asked from the top. The
first ``yes'' gives the answer, and the questions below it are not asked.}
\end{omfigure}

Because the first true \omterm{def:g4:algorithms-repeat-and-decide:condition}{condition} wins, the order of the cases matters. Here
the two questions are swapped: every child is also under 65, so the second
\omterm{def:g6:if-else:if}{branch} can never be reached.

\Needspace*{11\baselineskip}
\omcode{code/grade-6/04-if-else/listings/wrong_order.py}{1}{8}{The cases in the
wrong order.}\label{lst:g6:if-else:wrong-order}

\Needspace*{7\baselineskip}
\omoutput{code/grade-6/04-if-else/outputs/wrong_order.txt}{The child of 8 pays
the adult price.}\label{out:g6:if-else:wrong-order}

\section{\texorpdfstring{\texttt{and}}{and}, \texorpdfstring{\texttt{or}}{or}, \texorpdfstring{\texttt{not}}{not}}

\begin{definition}[Compound condition]\label{def:g6:if-else:compound}
A \emph{compound condition}\index{compound condition} joins \omterm{def:g4:algorithms-repeat-and-decide:condition}{conditions} with
\code{and}, \code{or} and \code{not}. \code{a and b} is true when both are
true; \code{a or b} is true when at least one of them is true; \code{not a} is
true when \code{a} is false.
\end{definition}

The \omterm{def:g4:programs-and-operating-system:program}{program} below prints, for each of the four ways to choose \code{a} and
\code{b}, the \omterm{def:g6:variables-and-input:value}{values} of \code{a and b} and of \code{a or b}.

\Needspace*{6\baselineskip}
\omcode{code/grade-6/04-if-else/listings/truth.py}{1}{3}{The four
cases.}\label{lst:g6:if-else:truth}

\Needspace*{8\baselineskip}
\omoutput{code/grade-6/04-if-else/outputs/truth.txt}{Columns: \code{a},
\code{b}, \code{a and b}, \code{a or b}.}\label{out:g6:if-else:truth}

\begin{omfigure}
\begin{tikzpicture}[x=1cm, y=1cm, font=\footnotesize]
  \foreach \t/\x/\v in {{\code{and}}/0/{0,0,0,1}, {\code{or}}/5/{0,1,1,1}} {
    \begin{scope}[shift={(\x,0)}]
      \node at (1.5,2.6) {\t};
      \node at (0.5,2.0) {\code{a}}; \node at (1.5,2.0) {\code{b}}; \node at (3.0,2.0) {result};
      \foreach \a/\b [count=\r] in {F/F, F/T, T/F, T/T} {
        \node at (0.5,2.0-0.5*\r) {\a}; \node at (1.5,2.0-0.5*\r) {\b};
      }
      \foreach \res [count=\r] in \v {
        \ifnum\res=1 \node[fill=ocMint!35, minimum width=1.2cm] at (3.0,2.0-0.5*\r) {T};
        \else \node[fill=black!8, minimum width=1.2cm] at (3.0,2.0-0.5*\r) {F}; \fi
      }
      \draw[black!50] (0,1.75) -- (3.7,1.75);
      \draw[black!50] (2.2,2.25) -- (2.2,-0.25);
    \end{scope}
  }
\end{tikzpicture}

{\small The four cases (T for \code{True}, F for \code{False}). \code{and}
needs both; \code{or} needs at least one.}
\end{omfigure}

A real rule that needs them is the rule of leap years: a year has a 29th of
February if it is divisible by 4, except the years that end a century, which
are leap years only when divisible by 400.
% ledger: g6:gregorian-rule

\Needspace*{8\baselineskip}
\omcode{code/grade-6/04-if-else/listings/leap.py}{1}{5}{Leap year or
not?}\label{lst:g6:if-else:leap}

\Needspace*{6\baselineskip}
\omoutput{code/grade-6/04-if-else/outputs/leap.txt}{1900 is divisible by 4,
but it ends a century and is not divisible by
400.}\label{out:g6:if-else:leap}

The brackets matter: Python works out \code{and} before \code{or}, like
$\times$ before $+$, so without them the \omterm{def:g4:algorithms-repeat-and-decide:condition}{condition} would mean something else.

\section{Testing every branch}

A \omterm{def:g4:programs-and-operating-system:program}{program} with \omterm{def:g6:if-else:if}{branches} can be right for one \omterm{def:g1:what-is-a-computer:input}{input} and wrong for another. To
test it, we choose \omterm{def:g1:what-is-a-computer:input}{inputs} that reach every \omterm{def:g6:if-else:if}{branch}, and we try the
\emph{boundaries}: the \omterm{def:g6:variables-and-input:value}{values} where the answer changes. For the cinema, they
are 12 and 65, and the ages just before them.

\Needspace*{12\baselineskip}
\omcode{code/grade-6/04-if-else/listings/boundaries.py}{1}{9}{Both sides of
each boundary.}\label{lst:g6:if-else:boundaries}

\Needspace*{8\baselineskip}
\omoutput{code/grade-6/04-if-else/outputs/boundaries.txt}{At 12 the child
price stops; at 65 the senior price
starts.}\label{out:g6:if-else:boundaries}

\begin{omfigure}
\begin{tikzpicture}[x=0.13cm, y=1cm, font=\footnotesize]
  \fill[ocBlue!25] (0,0) rectangle (12,0.4);
  \fill[ocYellow!45] (12,0) rectangle (65,0.4);
  \fill[ocMint!35] (65,0) rectangle (90,0.4);
  \draw[-{Stealth[length=2mm]}] (0,0) -- (92,0);
  \foreach \a in {0,12,65} \draw (\a,0) -- (\a,-0.15) node[below] {\a};
  \node at (6,0.75) {5 euros}; \node at (38.5,0.75) {9 euros}; \node at (77.5,0.75) {6 euros};
  \foreach \a in {11,12,64,65} \fill[ocRed] (\a,0.2) circle (0.07);
  \node[anchor=west] at (93,0) {age};
\end{tikzpicture}

{\small The price bands. The red dots are the four ages tried above, on both
sides of each boundary. Age 12 is in the middle band: \code{age < 12} is
false for it.}
\end{omfigure}

With a \omterm{def:g4:algorithms-repeat-and-decide:loop}{loop}, the whole rule can be checked at once:

\Needspace*{12\baselineskip}
\omcode{code/grade-6/04-if-else/listings/all_prices.py}{1}{9}{Ages 0 to 80, by
fives.}\label{lst:g6:if-else:all-prices}

\omoutput[1-5,13-18]{code/grade-6/04-if-else/outputs/all_prices.txt}{Each of the
three \omterm{def:g6:if-else:if}{branches} is reached.}\label{out:g6:if-else:all-prices}

\section{Exercises}

\begin{exercise}[$\star$]\label{exo:g6:if-else:1}
What do \code{print(3 < 5)} and \code{print(4 == 4.0)} print?
\end{exercise}

\begin{exercise}[$\star$]\label{exo:g6:if-else:2}
Which line does this \omterm{def:g4:programs-and-operating-system:program}{program} print?

\omcode{code/grade-6/04-if-else/exercises/which.py}{1}{7}{Which
\omterm{def:g6:if-else:if}{branch}?}\label{lst:g6:if-else:which}
\end{exercise}

\begin{exercise}[$\star$]\label{exo:g6:if-else:3}
With \code{a = True} and \code{b = False}, what is \code{not (a and b)}? And
\code{not a or b}?
\end{exercise}

\begin{exercise}[$\star$]\label{exo:g6:if-else:4}
Write a \omterm{def:g4:programs-and-operating-system:program}{program} that asks for a number and prints \code{positive or zero} or
\code{negative}.
\end{exercise}

\begin{exercise}[$\star\star$]\label{exo:g6:if-else:5}
Lina writes \code{if x = 3:} and gets a \code{SyntaxError}. Why? Fix it.
\end{exercise}

\begin{exercise}[$\star\star$]\label{exo:g6:if-else:6}
This \omterm{def:g4:programs-and-operating-system:program}{program} never prints \code{senior}, whatever the age. Why? Fix it.

\omcode{code/grade-6/04-if-else/exercises/senior_bug.py}{1}{7}{Where is the
senior?}\label{lst:g6:if-else:senior-bug}
\end{exercise}

\begin{exercise}[$\star\star$]\label{exo:g6:if-else:7}
Write a \omterm{def:g4:programs-and-operating-system:program}{program} that asks for three numbers and prints the largest.
\end{exercise}

\begin{exercise}[$\star\star$]\label{exo:g6:if-else:8}
Write a \omterm{def:g4:programs-and-operating-system:program}{program} that turns a mark out of 20 into a letter: A from 16, B from 12,
C from 8, D below 8.
\end{exercise}

\begin{exercise}[$\star\star$]\label{exo:g6:if-else:9}
Give a list of temperatures that reaches every \omterm{def:g6:if-else:if}{branch} of this \omterm{def:g4:programs-and-operating-system:program}{program}, with
the boundaries.

\omcode{code/grade-6/04-if-else/exercises/reach.py}{1}{9}{Four
\omterm{def:g6:if-else:if}{branches}.}\label{lst:g6:if-else:reach}
\end{exercise}

\begin{exercise}[$\star\star\star$]\label{exo:g6:if-else:10}
A \omterm{def:g5:networks-first-look:password}{password} should have at least 12 characters, and \code{len(password)} gives
the number of characters of a text. Write a \omterm{def:g4:programs-and-operating-system:program}{program} that asks for a \omterm{def:g5:networks-first-look:password}{password}
and says whether it is long enough, and if not, how many characters are
missing.
\end{exercise}

\begin{exercise}[$\star\star\star$]\label{exo:g6:if-else:11}
The \omterm{def:g4:programs-and-operating-system:program}{program} below decides ``yes'' or ``no'' in two ways: with an \code{if}
inside an \code{if}, and with \code{and}. Do the two always agree? Check it on
all four cases, by hand and by running it.

\omcode{code/grade-6/04-if-else/exercises/nested_vs_and.py}{1}{14}{Two ways
to decide.}\label{lst:g6:if-else:nested-vs-and}
\end{exercise}

\begin{exercise}[$\star\star\star$]\label{exo:g6:if-else:12}
Tom's leap-year \omterm{def:g4:programs-and-operating-system:program}{program} says that 1900 is a leap year. Find the mistake, fix
it, and give a year that shows the fix is right.

\omcode{code/grade-6/04-if-else/exercises/leap_bug.py}{1}{5}{Wrong for
1900.}\label{lst:g6:if-else:leap-bug}
\end{exercise}

\section{Problem: The Calendar's Leap Years}

\begin{problem}[{Weekend problem --- 97 leap years}]\label{pb:g6:if-else:1}
The Earth goes round the Sun in a little less than 365 days and a quarter. A
calendar must give each year a whole number of days, and the leap years add a
day from time to time so that the seasons stay in the same months. Since 1582,
most of the world uses the rule of \cref{lst:g6:if-else:leap}.
% ledger: g6:gregorian-1582

\medskip\noindent\textbf{Part I --- By hand.}
\begin{enumerate}
  \item Is 2024 a leap year? Is 2023?
  \item Is 2100 a leap year? Is 2000? Is 1900? Give the reason each time.
  \item How many days does a leap year have? How many do four years in a row
    have, if exactly one of them is a leap year?
  \item In the \omterm{def:g4:algorithms-repeat-and-decide:condition}{condition} of \cref{lst:g6:if-else:leap}, which part is true
    for 2000, and which for 2024?
\end{enumerate}

\medskip\noindent\textbf{Part II --- A program.}
\begin{enumerate}[resume]
  \item Run \cref{lst:g6:if-else:leap} for your four years of question 2. Do
    the answers agree with yours?
  \item Write a \omterm{def:g4:programs-and-operating-system:program}{program} that counts the leap years from 1901 to 2100, both
    included, with a \code{for} \omterm{def:g4:algorithms-repeat-and-decide:loop}{loop}, an \code{if} and a \omterm{def:g6:variables-and-input:value}{variable}
    \code{count} that grows by 1 for each leap year.
  \item Change it to count from 2001 to 2400. Why are there more?
  \item Explain why every block of 400 years in a row, such as 2001 to 2400 or
    1601 to 2000, has the same number of leap years.
\end{enumerate}

\medskip\noindent\textbf{Part III --- The average year.}
\begin{enumerate}[resume]
  \item Count the leap years in 400 years by hand: how many years are
    divisible by 4, how many of them end a century, and how many of those are
    divisible by 400?
  \item How many days are there in 400 years?
  \item What is the average length of a year, in days? Compare it with $365
    + 1/4$.
  \item State the result.
\end{enumerate}
\end{problem}
